% !PRINTED 204
\weberpage{204}
%<<CONT
\weq{4}{\begin{gathered}
H + 4\,a_2'\,f = -3\,\psi_1^2\\
H + 4\,a_2''\,f = -3\,\psi_2^2\\
H + 4\,a_2'''\,f = -3\,\psi_3^2.
\end{gathered}}

Von den drei Functionen $\psi_1, \psi_2, \psi_3$ können \gesperrt{keine zwei
einen gemeinschaftlichen Theiler haben}, wenn die Discriminante
$D$ von Null verschieden vorausgesetzt wird. Denn
dann sind auch die drei Grössen $a_2', a_2'', a_2'''$ von einander verschieden
. Wenn also $\psi_1$ und $\psi_2$ verschwinden, so müssen $H$ und
$f$ verschwinden, d.\,h. ein gemeinsamer Theiler von $\psi_1$ und $\psi_2$
müsste gemeinsamer Theiler von $H$ und $f$ sein. Nun war aber $\xi$
ein beliebiger Lineartheiler von $f$. Dieser kann nach (2) nur
dann Theiler von $H$ sein, wenn $a_3' = 0$ ist, also wieder, wenn
die Discriminante $D$ verschwindet [§.\,65, (11)].

Nun ist die Functionaldeterminante $T$ von $f$ und $H$ identisch
mit der Functionaldeterminante von $f$ und $H + \lambda f$, was auch
$\lambda$ sein mag, und daraus ergiebt sich, wenn $\lambda = 4\,a_2'$ gesetzt
wird,
\[
T = -\tfrac{1}{2}\,\psi_1\,[F'(\xi)\,\psi_1'(\eta) - F'(\eta)\,\psi_1'(\xi)].
\]

Es ist also $T$ theilbar durch $\psi_1$, und aus gleichen Gründen
durch $\psi_2$ und $\psi_3$, und folglich ist $T$ bis auf einen von $\xi, \eta$ unabhängigen
 Factor identisch mit dem Product $\psi_1\,\psi_2\,\psi_3$.
Bilden wir demnach das Product der drei Gleichungen (4),
so folgt mit Rücksicht auf die cubische Gleichung §.\,65, (12),
der die Grössen $a_2', a_2'', a_2'''$ genügen:
\[
H^3 - 48\,A\,H\,f^2 - 64\,B\,f^3 = c\,T^2,
\]
wo $c$ eine noch zu bestimmende Constante ist. Nach §.\,64, (12)
bleibt $c$ bei einer linearen Transformation ungeändert, und man
findet also seinen Werth, wenn man aus der Normalform (1), (2)
die Glieder mit der höchsten Potenz von $\xi$ beiderseits einander
gleich setzt, $c = -27$.

Wir haben also zwischen den Invarianten und Covarianten
die folgende identische Relation
\weq{5}{H^3 - 48\,A\,H\,f^2 - 64\,B\,f^3 = -27\,T^2.}

Die Functionen $\psi_1, \psi_2, \psi_3$ sind gleichfalls Covarianten, freilich
 mit Coëfficienten, die in den Coëfficienten von $f$ nicht rational
 sind. Wir können sie leicht durch die Wurzeln $\alpha, \beta, \gamma, \delta$
von $f = 0$ ausdrücken. Setzt man zur Vereinfachung $y = 1$,
so folgt nach §.\,65, (3) und (5)
%CONT>>
